% Lecture 9 notes: \input by notes.tex, \subimport'ed by the semester book.
% Exam questions (экзамен): 21 (быстрое возведение точки в степень, тернарное разложение,
% оценки числа операций), 22 (ECDLP, тривиальная оценка, ECDH), 23 (rho-алгоритм для ECDLP).
\begin{lecture}{9}{Возведение точки в степень и задача ECDLP}

\section{Зачем это нужно}

В лекции~8 мы построили группу $E(\FF_p)$ точек эллиптической кривой
$E\colon y^2 = x^3 + ax + b$ над конечным полем: сумма двух точек вычисляется по явным формулам
через несколько умножений и одно обращение в $\FF_p$, а по теореме Хассе
$\abs{\#E(\FF_p) - p - 1} \le 2\sqrt p$. Группа $E(\FF_p)$ коммутативна, и операцию в ней
принято записывать \emph{аддитивно}: вместо $g^n$ пишут
\[
  nP = \underbrace{P + P + \cdots + P}_{n}, \qquad 0 \cdot P = \ECpt, \qquad (-n)P = n(-P).
\]
Мы по-прежнему будем говорить о «возведении точки в степень» (англ.\ \emph{scalar
multiplication}).

Все протоколы на эллиптических кривых устроены так же, как протоколы в $\Znx{p}$ (лекция~1):
\begin{itemize}
  \item «законные» участники вычисляют $nP$ по известным $n$ и $P$ "--- это должно быть
    \emph{быстро}: за $\bigO(\log n)$ групповых операций;
  \item противник, видя $P$ и $Q = nP$, должен найти $n$ "--- это \emph{задача
    дискретного логарифмирования на эллиптической кривой} (ECDLP), и она должна быть
    \emph{трудной}.
\end{itemize}
Главная причина популярности эллиптических кривых (Н.~Коблиц~\cite{koblitz-ecc} и
В.~Миллер~\cite{miller-ecc}, 1985) в том, что для ECDLP на «хорошей» кривой не известно
алгоритмов быстрее универсальных, работающих за $\bigO(\sqrt N)$ операций, где $N$ "--- порядок
точки. Для $\Znx{p}$ есть субэкспоненциальный метод исчисления индексов (лекция~4), поэтому
при одинаковой стойкости ключи ECC в несколько раз короче.

В этой лекции: алгоритм «удвоение"--~сложение» и его оценки (раздел~\ref{l9:dbladd}),
тернарное разложение показателя и несмежная форма (раздел~\ref{l9:naf}), задача ECDLP и
протоколы ECDH, ECElGamal, ECDSA (разделы~\ref{l9:ecdlp}--\ref{l9:protocols}), $\rho$-метод
Полларда для ECDLP (раздел~\ref{l9:rho}).

\section{Алгоритм «удвоение"--~сложение»}\label{l9:dbladd}

Пусть $n = (n_{L-1} \dots n_1 n_0)_2$, $n_{L-1} = 1$, т.\,е. $L = \floor{\log_2 n} + 1$ "---
двоичная длина $n$. Как и бинарный алгоритм возведения в степень (лекция~6 первого семестра),
вычисление $nP$ опирается на схему Горнера
\[
  nP = 2\bigl(\cdots 2\bigl(2P + n_{L-2}P\bigr) + n_{L-3}P \cdots\bigr) + n_0 P.
\]

\begin{algorithm}[«удвоение"--~сложение», слева направо]\label{l9:alg-dbladd}\leavevmode
\begin{algorithmic}[1]
\Require $P \in E(\FF_p)$, $n = (n_{L-1} \dots n_0)_2 \ge 1$.
\Ensure $R = nP$.
\State $R \gets P$
\For{$i = L - 2, L - 3, \dots, 0$}
  \State $R \gets 2R$ \Comment{удвоение}
  \If{$n_i = 1$} \State $R \gets R + P$ \Comment{сложение} \EndIf
\EndFor
\State \Return $R$
\end{algorithmic}
\end{algorithm}

\begin{proposition}\label{l9:dbladd-correct}
Алгоритм~\ref{l9:alg-dbladd} вычисляет $nP$.
\end{proposition}

\begin{proof}
Пусть $m_i = (n_{L-1} \dots n_i)_2$ "--- число, составленное из старших $L - i$ битов. Индукцией
по $i = L - 1, L - 2, \dots, 0$ проверим, что после обработки бита $n_i$ выполнено
$R = m_i P$. При $i = L - 1$: $R = P = m_{L-1} P$. Шаг: $m_i = 2m_{i+1} + n_i$, а алгоритм
выполняет $R \gets 2R$ и при $n_i = 1$ прибавляет $P$. В конце $R = m_0 P = nP$.
\end{proof}

Удвоение и сложение точек вычисляются по разным формулам (лекция~8):
\[
\begin{gathered}
  \lambda = \frac{3x_1^2 + a}{2y_1}\ \text{(удвоение)}, \qquad
  \lambda = \frac{y_2 - y_1}{x_2 - x_1}\ \text{(сложение)},\\
  x_3 = \lambda^2 - x_1 - x_2, \qquad y_3 = \lambda(x_1 - x_3) - y_1,
\end{gathered}
\]
но по стоимости они близки (одно обращение и 2--3 умножения в $\FF_p$), поэтому мы считаем
обе операции \emph{элементарными групповыми операциями}.

\begin{theorem}\label{l9:dbladd-cost}
Пусть $n \ge 1$, $L = \floor{\log_2 n} + 1$, $w(n)$ "--- число единиц в двоичной записи $n$
(вес Хэмминга). Алгоритм~\ref{l9:alg-dbladd} выполняет $L - 1$ удвоений и $w(n) - 1$ сложений.
Следовательно:
\begin{enumerate}
  \item максимальное число групповых операций равно $2(L - 1) = 2\floor{\log_2 n}$
    (достигается при $n = 2^L - 1$);
  \item если $n$ выбирается случайно и равновероятно среди $L$-битовых чисел
    ($2^{L-1} \le n < 2^L$), то математическое ожидание числа операций равно
    $\frac32 (L - 1) \approx 1{,}5 \log_2 n$.
\end{enumerate}
\end{theorem}

\begin{proof}
Цикл выполняется $L - 1$ раз, на каждой итерации одно удвоение и, если $n_i = 1$, одно
сложение. Единиц среди $n_{L-2}, \dots, n_0$ ровно $w(n) - 1$ (старший бит учтён начальным
присваиванием $R \gets P$). Итого $T(n) = (L - 1) + (w(n) - 1) \le 2(L - 1)$, равенство при
$w(n) = L$, т.\,е. $n = 2^L - 1$.

Если $n$ равномерно распределено среди $L$-битовых чисел, то биты $n_0, \dots, n_{L-2}$
независимы и каждый равен $1$ с вероятностью $\frac12$. По линейности математического
ожидания $\mathbb{E}\,(w(n) - 1) = \sum_{i=0}^{L-2} \Pr[n_i = 1] = \frac{L-1}{2}$, откуда
$\mathbb{E}\,T(n) = (L - 1) + \frac{L - 1}{2} = \frac32 (L - 1)$.
\end{proof}

\begin{example}\label{l9:ex-15P}
Возьмём кривую $E\colon y^2 = x^3 + 2x + 2$ над $\FF_{17}$ (её мы будем использовать всю
лекцию). Здесь $\#E(\FF_{17}) = 19$ "--- простое число, поэтому любая точка $P \ne \ECpt$
имеет порядок $19$. Пусть $P = (5, 1)$, $n = 15 = 1111_2$. Все вычисления по модулю $17$:
\[
\begin{array}{l|r|r|c}
  \text{операция} & k & \lambda & R = kP \\ \hline
  2R    & 2  & 13 & (6, 3)   \\
  R + P & 3  & 2  & (10, 6)  \\
  2R    & 6  & 11 & (16, 13) \\
  R + P & 7  & 15 & (0, 6)   \\
  2R    & 14 & 3  & (9, 1)   \\
  R + P & 15 & 0  & (3, 16)
\end{array}
\]
Например, первое удвоение: $\lambda = (3 \cdot 5^2 + 2)/(2 \cdot 1) = 77 \cdot 2^{-1} \equiv 9
\cdot 9 = 81 \equiv 13$, $x_3 = 13^2 - 10 = 159 \equiv 6$, $y_3 = 13(5 - 6) - 1 = -14 \equiv 3$.
Итак, $15P = (3, 16)$ за $3 + 3 = 6$ операций "--- максимум для $L = 4$.
\end{example}

\taskexample{ecdlp-binary}

\begin{remark}
Существует и вариант «справа налево» (именно он приведён в~\cite[разд.~6.3.1]{hps}):
$Q \gets P$, $R \gets \ECpt$; пока $n > 0$: если $n$ нечётно, $R \gets R + Q$; затем
$Q \gets 2Q$, $n \gets \floor{n/2}$. Он выполняет столько же операций, но хранит две
переменные точки; вариант «слева направо» прибавляет всегда одну и ту же точку $P$, что удобно
для окон и предвычислений.
\end{remark}

\checkyourself{Сколько удвоений и сложений нужно для вычисления $1000P$ алгоритмом
«удвоение"--~сложение»? Почему в оценке среднего числа операций стоит $L - 1$, а не $L$?}

\section{Тернарное разложение и несмежная форма}\label{l9:naf}

\subsection{Вычитание точки бесплатно}

В $\Znx{p}$ обращение $g^{-1}$ стоит заметно дороже умножения, а на эллиптической кривой
противоположная точка вычисляется даром:
\[
  -P = -(x, y) = (x, -y),
\]
т.\,е. вычитание точки $R - P = R + (-P)$ стоит ровно столько же, сколько сложение. Поэтому
для кривых выгодно записывать показатель с цифрами $-1, 0, 1$: например,
$15 = 1111_2 = 2^4 - 1$, и вместо трёх сложений достаточно одного вычитания.

\begin{definition}
\emph{Тернарным} (знаковым двоичным) разложением целого $n \ge 1$ называется запись
\[
  n = \sum_{i=0}^{\ell - 1} z_i 2^i, \qquad z_i \in \{-1, 0, 1\}, \quad z_{\ell-1} \ne 0.
\]
Пишут $n = (z_{\ell-1} \dots z_0)$, обозначая $-1$ через $\bar 1$. Разложение называется
\emph{несмежной формой} (англ.\ \emph{non-adjacent form}, NAF), если из любых двух соседних
цифр хотя бы одна равна нулю: $z_i z_{i+1} = 0$ при всех $i$.
\end{definition}

Тернарное разложение неединственно: $3 = (11) = (10\bar 1)$, $7 = (111) = (100\bar 1) =
(1\bar 1 11)$. Несмежная форма выделяет среди них одно.

\begin{theorem}\label{l9:naf-exists}
Каждое натуральное $n$ имеет единственную несмежную форму. Её цифры вычисляются справа налево:
\[
  z_i = \begin{cases} 0, & n_i \text{ чётно},\\ 2 - (n_i \bmod 4) \in \{1, -1\}, & n_i
  \text{ нечётно},\end{cases} \qquad n_0 = n, \quad n_{i+1} = \frac{n_i - z_i}{2},
\]
пока $n_i \ne 0$.
\end{theorem}

\begin{proof}
\emph{Существование.} Числа $n_i$ целые: если $n_i$ нечётно, то $z_i \equiv n_i \pmod 2$, и
$n_i - z_i$ чётно. Последовательность конечна: при $n_i \ge 2$ имеем
$n_{i+1} \le (n_i + 1)/2 < n_i$, а $n_i = 1$ даёт $z_i = 1$, $n_{i+1} = 0$; все $n_i \ge 0$, так
как $n_i - z_i \ge 0$ при $n_i \ge 1$. По построению $n_i = z_i + 2n_{i+1}$, поэтому
$n = \sum_i z_i 2^i$. Несмежность: если $z_i \ne 0$, то $z_i \equiv n_i \pmod 4$
(при $n_i \equiv 1$ имеем $z_i = 1$, при $n_i \equiv 3$ "--- $z_i = -1 \equiv 3$), значит,
$4 \mid n_i - z_i$, число $n_{i+1}$ чётно и $z_{i+1} = 0$.

\emph{Единственность.} Пусть $n = \sum_{i \ge 0} z_i 2^i$ "--- несмежная форма. Докажем, что
$z_0$ определяется числом $n$ однозначно и совпадает с цифрой из алгоритма; тогда
$(n - z_0)/2 = \sum_{i \ge 1} z_i 2^{i-1}$ "--- несмежная форма меньшего числа, и
утверждение следует по индукции. Если $n$ чётно, то $z_0 \equiv n \equiv 0 \pmod 2$, т.\,е.
$z_0 = 0$. Если $n$ нечётно, то $z_0 = \pm 1$, а по несмежности $z_1 = 0$, поэтому
$n = z_0 + 4\sum_{i \ge 2} z_i 2^{i-2} \equiv z_0 \pmod 4$. Из двух чисел $\pm 1$ ровно одно
сравнимо с $n$ по модулю $4$: $z_0 = 1$ при $n \equiv 1$, $z_0 = -1$ при $n \equiv 3 \pmod 4$.
Это и есть $2 - (n \bmod 4)$.
\end{proof}

\begin{example}\label{l9:ex-naf}
$n = 1979 = 11110111011_2$:
\[
\begin{array}{c|rrrrrrrrrrrr}
  i   & 0 & 1 & 2 & 3 & 4 & 5 & 6 & 7 & 8 & 9 & 10 & 11 \\ \hline
  n_i & 1979 & 990 & 495 & 248 & 124 & 62 & 31 & 16 & 8 & 4 & 2 & 1 \\
  z_i & \bar 1 & 0 & \bar 1 & 0 & 0 & 0 & \bar 1 & 0 & 0 & 0 & 0 & 1
\end{array}
\]
т.\,е. $1979 = (10000\bar 1 000\bar 1 0\bar 1) = 2^{11} - 2^6 - 2^2 - 1$. В двоичной записи
$9$ единиц, а в NAF "--- $4$ ненулевые цифры, хотя NAF на одну цифру длиннее.
\end{example}

\taskexample{ecdlp-naf}

\begin{algorithm}[«удвоение"--~сложение"--~вычитание»]\label{l9:alg-naf}\leavevmode
\begin{algorithmic}[1]
\Require $P \in E(\FF_p)$, несмежная форма $n = (z_{\ell-1} \dots z_0)$, $z_{\ell-1} = 1$.
\Ensure $R = nP$.
\State $R \gets P$, \quad $\bar P \gets -P = (x_P, -y_P)$
\For{$i = \ell - 2, \dots, 0$}
  \State $R \gets 2R$
  \If{$z_i = 1$} \State $R \gets R + P$ \EndIf
  \If{$z_i = -1$} \State $R \gets R + \bar P$ \EndIf
\EndFor
\State \Return $R$
\end{algorithmic}
\end{algorithm}

Корректность доказывается так же, как в утверждении~\ref{l9:dbladd-correct}: после обработки
цифры $z_i$ выполнено $R = \bigl(\sum_{j \ge i} z_j 2^{j - i}\bigr) P$.

\begin{example}\label{l9:ex-15P-naf}
На той же кривой $y^2 = x^3 + 2x + 2$ над $\FF_{17}$, $P = (5, 1)$: $15 = (1000\bar 1)$,
$-P = (5, 16)$.
\[
\begin{array}{l|r|r|c}
  \text{операция} & k & \lambda & R = kP \\ \hline
  2R    & 2  & 13 & (6, 3)   \\
  2R    & 4  & 7  & (3, 1)   \\
  2R    & 8  & 6  & (13, 7)  \\
  2R    & 16 & 6  & (10, 11) \\
  R - P & 15 & 16 & (3, 16)
\end{array}
\]
Получили тот же ответ $15P = (3, 16)$ за $5$ операций вместо $6$.
\end{example}

\taskexample{ecdlp-ternary}

\subsection{Оценки числа операций}

\begin{lemma}\label{l9:naf-len}
Пусть $n \ge 1$ имеет двоичную длину $L$, а его несмежная форма "--- длину $\ell$. Тогда
$\ell \le L + 1$, и в NAF не более $\ceil{\ell/2}$ ненулевых цифр.
\end{lemma}

\begin{proof}
Старшая цифра NAF положительного числа равна $1$ (иначе сумма была бы отрицательной, см.\
оценку ниже), а $z_{\ell-2} = 0$. По несмежности младшие цифры по модулю дают не больше
\[
  \sum_{t \ge 0} 2^{\ell - 3 - 2t} < 2^{\ell - 3} \cdot \frac{1}{1 - 1/4} = \frac13 \, 2^{\ell-1}.
\]
Поэтому $n > 2^{\ell-1} - \frac13 2^{\ell-1} = \frac23 2^{\ell-1}$, т.\,е.
$2^{\ell-1} < \frac32 n < \frac32 \cdot 2^L < 2^{L+1}$, откуда $\ell - 1 < L + 1$ и $\ell \le L + 1$.
(Та же оценка с $z_{\ell-1} = -1$ дала бы $n < 0$.) Ненулевые цифры занимают позиции, попарно
отстоящие хотя бы на $2$, а позиций всего $\ell$; значит, их не больше $\ceil{\ell/2}$.
\end{proof}

Например, $n = 2^L - 1 = (1 0 \dots 0 \bar 1)$: NAF длиннее двоичной записи на одну цифру.
Существенно, что NAF \emph{никогда} не содержит больше ненулевых цифр, чем любое другое тернарное
разложение того же числа (минимальность веса NAF; доказательство "--- см.~\cite[разд.~14.7]{hac},
\cite[гл.~11]{galbraith}).

Для оценки среднего веса NAF нам понадобится наблюдение о «равномерности» чисел $n_i$.

\begin{lemma}\label{l9:naf-uniform}
Пусть $n$ равномерно распределено на множестве $\{0, 1, \dots, 2^m - 1\}$, а цифры
$z_0, \dots, z_{k-1}$, $k < m$, получены алгоритмом теоремы~\ref{l9:naf-exists} (при $n_i = 0$
считаем все дальнейшие цифры нулевыми). Если $k = 0$ или $z_{k-1} = 0$, то при любых
фиксированных значениях $z_0, \dots, z_{k-1}$ число $n_k$ нечётно с вероятностью ровно $\frac12$.
\end{lemma}

\begin{proof}
Выбор цифры $z_i$ определяется вычетом $n_i \bmod 4$, а $n_i = (n - c_i)/2^i$, где
$c_i = \sum_{j < i} z_j 2^j$; по индукции $z_0, \dots, z_{k-2}$ определяются вычетом
$n \bmod 2^k$, а при $z_{k-1} = 0$ "--- лишь чётностью $n_{k-1}$, т.\,е. тоже вычетом
$n \bmod 2^k$. Следовательно, множество тех $n$, для которых цифры равны данным, "--- это
\emph{весь} класс вычетов $n \equiv c_k \pmod{2^k}$ (если не пусто). В нём $2^{m-k}$ чисел
$n = c_k + 2^k t$, и $n_k = t$ пробегает $2^{m-k}$ подряд идущих целых значений (с точностью
до сдвига), т.\,е. ровно половина из них нечётна.
\end{proof}

\begin{theorem}\label{l9:naf-density}
Пусть $n$ равномерно распределено на $\{0, \dots, 2^m - 1\}$. Тогда математическое ожидание
числа ненулевых цифр NAF числа $n$ равно $\frac m3 + \bigO(1)$: средняя плотность ненулевых
цифр NAF равна $\frac13$ (против $\frac12$ у двоичной записи).
\end{theorem}

\begin{proof}
Назовём позицию $k$ \emph{свободной}, если $k = 0$ или $z_{k-1} = 0$; иначе (после ненулевой
цифры) $z_k = 0$ по несмежности. По лемме~\ref{l9:naf-uniform} в свободной позиции $z_k \ne 0$
с вероятностью $\frac12$ независимо от предыдущих цифр. Пусть $q_k$ "--- вероятность того,
что позиция $k$ свободна. Тогда $q_0 = 1$ и позиция $k + 1$ свободна, если позиция $k$
несвободна или свободна и $z_k = 0$:
\[
  q_{k+1} = (1 - q_k) + \tfrac12 q_k = 1 - \tfrac12 q_k, \qquad\text{откуда}\qquad
  q_k - \tfrac23 = \left(-\tfrac12\right)^k \left(q_0 - \tfrac23\right) = \tfrac13 \left(-\tfrac12\right)^k.
\]
Ожидаемое число ненулевых цифр в позициях $0, \dots, m - 1$ равно
\[
  \sum_{k=0}^{m-1} \Pr[z_k \ne 0] = \sum_{k=0}^{m-1} \tfrac12 q_k
  = \frac m3 + \frac16 \sum_{k=0}^{m-1} \left(-\tfrac12\right)^k = \frac m3 + \bigO(1),
\]
и ещё не более одной ненулевой цифры может стоять в позиции $m$ (лемма~\ref{l9:naf-len}).
\end{proof}

Наглядно: справа налево NAF разбивается на блоки «$0$» и «$0 z$», $z = \pm 1$ (ненулевая
цифра вместе с обязательным нулём слева от неё); каждый блок появляется с вероятностью
$\frac12$, средняя длина блока $\frac32$, а ненулевых цифр в блоке в среднем $\frac12$ "---
отсюда плотность $\frac12 : \frac32 = \frac13$. Проверка на компьютере: для всех
$n < 2^{16}$ среднее число ненулевых цифр NAF равно $5{,}78 \approx 16/3 + 0{,}44$.

\begin{corollary}\label{l9:naf-cost}
Вычисление $nP$ по NAF (алгоритм~\ref{l9:alg-naf}) требует $\ell - 1 \le \log_2 n + 1$ удвоений
и (число ненулевых цифр $- 1$) сложений/вычитаний. Поэтому
\begin{enumerate}
  \item максимальное число групповых операций не превосходит $\frac32 \log_2 n + \bigO(1)$;
  \item среднее число групповых операций равно $\frac43 \log_2 n + \bigO(1)$.
\end{enumerate}
Для сравнения, у бинарного метода $2\log_2 n$ и $\frac32 \log_2 n$ соответственно
(теорема~\ref{l9:dbladd-cost}).
\end{corollary}

\begin{proof}
По лемме~\ref{l9:naf-len} $\ell \le L + 1 \le \log_2 n + 2$, и ненулевых цифр не больше
$\ceil{\ell/2} \le \frac12 \log_2 n + \frac32$; суммируя, получаем первую оценку. Для второй
заметим, что цифра $z_k$ зависит лишь от $n \bmod 2^{k+2}$. Если $n$ равномерно распределено
среди $L$-битовых чисел, то $n' = n \bmod 2^{L-1}$ равномерно распределено на
$\{0, \dots, 2^{L-1} - 1\}$, и цифры $z_0, \dots, z_{L-3}$ у $n$ и $n'$ совпадают. По
теореме~\ref{l9:naf-density} с $m = L - 1$ среди них в среднем $\frac L3 + \bigO(1)$ ненулевых,
а оставшиеся не более трёх старших цифр дают $\bigO(1)$. Итого в среднем
$\frac L3 + \bigO(1)$ сложений и $L + \bigO(1)$ удвоений, т.\,е. $\frac43 \log_2 n + \bigO(1)$
операций.
\end{proof}

\begin{remark}[окна]
Обобщение "--- \emph{оконная} несмежная форма $w$-NAF: цифры "--- нечётные числа
$\abs{z_i} < 2^{w-1}$, среди любых $w$ подряд идущих цифр не более одной ненулевой. Средняя
плотность $\frac1{w+1}$, но заранее вычисляются точки $3P, 5P, \dots, (2^{w-1} - 1)P$.
\end{remark}

\begin{remark}[лесенка Монтгомери]
Время работы алгоритмов~\ref{l9:alg-dbladd} и~\ref{l9:alg-naf} зависит от битов секретного
показателя: по времени или энергопотреблению можно восстановить ключ (атаки по побочным
каналам, П.~Кохер~\cite{kocher-timing}). \emph{Лесенка Монтгомери}~\cite{montgomery-ladder}
хранит пару $(R_0, R_1)$ с инвариантом $R_1 - R_0 = P$:
\[
  R_0 \gets \ECpt,\ R_1 \gets P; \qquad\text{для } i = L - 1, \dots, 0:\quad
  \begin{cases} R_1 \gets R_0 + R_1,\ R_0 \gets 2R_0, & n_i = 0,\\
  R_0 \gets R_0 + R_1,\ R_1 \gets 2R_1, & n_i = 1, \end{cases}
\]
и в конце $R_0 = nP$. На каждый бит "--- ровно одно сложение и одно удвоение
(итого $2\log_2 n$ операций, как худший случай бинарного метода), зато последовательность
операций не зависит от битов~\cite{joye-yen-ladder}. На кривых Монтгомери
$By^2 = x^3 + Ax^2 + x$ лесенку можно вести по одной координате $x$ (так устроен протокол
X25519).
\end{remark}

\begin{remark}[проективные координаты]
Каждое сложение и удвоение в аффинных координатах требует обращения в $\FF_p$, а обращение
(алгоритм Евклида) в десятки раз дороже умножения. Поэтому на практике точки хранят в
\emph{проективных} координатах $(X : Y : Z)$, $x = X/Z$, $y = Y/Z$, или в \emph{якобиевых},
$x = X/Z^2$, $y = Y/Z^3$: формулы сложения и удвоения становятся формулами без деления
(около десятка умножений), а единственное обращение выполняется в конце, при переходе
к аффинной точке. Точка $\ECpt = (0 : 1 : 0)$ при этом "--- обычная точка, а не особый случай.
\end{remark}

\checkyourself{Найдите NAF чисел $7$, $11$ и $2^{10} - 1$. Сколько операций экономит NAF при
вычислении $(2^{10} - 1)P$?}

\checkyourself{Почему для возведения в степень в $\Znx{p}$ тернарное разложение даёт меньше
выгоды, чем для эллиптических кривых?}

\section{Задача ECDLP}\label{l9:ecdlp}

\begin{definition}
Пусть $E$ "--- эллиптическая кривая над $\FF_p$, $P \in E(\FF_p)$, $Q \in \langle P \rangle$.
\emph{Задача дискретного логарифмирования на эллиптической кривой} (ECDLP): найти целое $n$,
для которого $Q = nP$. Число $n$ определено по модулю $N = \ord P$ и обозначается
$\log_P Q$.
\end{definition}

Так как $nP = mP \iff N \mid n - m$, отображение $\Zn{N} \to \langle P \rangle$,
$n \mapsto nP$, "--- изоморфизм групп, и $\log_P(Q_1 + Q_2) \equiv \log_P Q_1 + \log_P Q_2
\pmod N$.

\begin{proposition}[тривиальная оценка]\label{l9:trivial}
ECDLP в группе $\langle P \rangle$ порядка $N$ решается полным перебором за не более чем $N - 1$
групповых операций и $\bigO(1)$ памяти. Так как $N \le \#E(\FF_p) \le p + 1 + 2\sqrt p$,
это $\bigO(p)$ операций.
\end{proposition}

\begin{proof}
Вычисляем $P, 2P = P + P, 3P = 2P + P, \dots$, сравнивая каждую точку с $Q$; так как
$Q \in \langle P \rangle = \{\ECpt, P, \dots, (N - 1)P\}$, совпадение наступит не позже чем
через $N - 1$ сложений (случай $Q = \ECpt$ проверяется сразу). Оценка $N$ "--- теорема Хассе
(лекция~8).
\end{proof}

\begin{example}\label{l9:ex-trivial}
На кривой $y^2 = x^3 + 2x + 2$ над $\FF_{17}$, $P = (5, 1)$ (порядок $19$), кратные $kP$:
\[
\begin{array}{c|ccccccccc}
  k  & 1 & 2 & 3 & 4 & 5 & 6 & 7 & 8 & 9 \\ \hline
  kP & (5,1) & (6,3) & (10,6) & (3,1) & (9,16) & (16,13) & (0,6) & (13,7) & (7,6)
\end{array}
\]
\[
\begin{array}{c|ccccccccc}
  k  & 10 & 11 & 12 & 13 & 14 & 15 & 16 & 17 & 18 \\ \hline
  kP & (7,11) & (13,10) & (0,11) & (16,4) & (9,1) & (3,16) & (10,11) & (6,14) & (5,16)
\end{array}
\]
и $19P = \ECpt$. Видно, что $(19 - k)P = -kP$. Например, $\log_P (16, 4) = 13$.
\end{example}

\subsection{Универсальные алгоритмы и выбор кривой}

Все \emph{универсальные} (англ.\ \emph{generic}) алгоритмы дискретного логарифмирования из
лекций~1--3 используют лишь групповую операцию и сравнение элементов, поэтому дословно
переносятся на $E(\FF_p)$:
\begin{itemize}
  \item метод Шенкса (малые и большие шаги) "--- $\bigO(\sqrt N)$ операций и
    $\bigO(\sqrt N)$ памяти;
  \item $\rho$-метод Полларда "--- $\bigO(\sqrt N)$ операций и $\bigO(1)$ памяти
    (раздел~\ref{l9:rho});
  \item алгоритм Полига"--~Хеллмана сводит ECDLP в группе порядка $N = \prod q_i^{e_i}$
    к задачам в группах порядков $q_i$: $\bigO\bigl(\sum e_i (\log N + \sqrt{q_i})\bigr)$ операций.
\end{itemize}
Для универсальных алгоритмов это неулучшаемо: В.~Шоуп~\cite{shoup-generic} доказал, что любой
универсальный алгоритм требует $\Omega(\sqrt q)$ групповых операций, где $q$ "--- наибольший
простой делитель порядка группы (без доказательства).

Отсюда правило выбора параметров: \emph{порядок базовой точки должен быть большим простым
числом}. Обычно $\#E(\FF_p) = hq$, где $q$ "--- простое порядка $p$, а \emph{кофактор}
$h \in \{1, 2, 4, 8\}$, и $P$ выбирают порядка $q$. Например, у стандартной кривой P-256
$h = 1$, у Curve25519 $h = 8$. Чтобы получить стойкость $2^{128}$ (столько операций нужно
противнику), достаточно $q \approx 2^{256}$, т.\,е. $p$ длиной $256$ бит.

\begin{table}[h]
\centering
\begin{tabular}{c|c|c}
  стойкость (бит) & RSA, DH в $\Znx{p}$: длина модуля & ECC: длина $q$ \\ \hline
  $80$  & $1024$  & $160$ \\
  $112$ & $2048$  & $224$ \\
  $128$ & $3072$  & $256$ \\
  $192$ & $7680$  & $384$ \\
  $256$ & $15360$ & $512$
\end{tabular}
\caption{Длины ключей одинаковой стойкости по рекомендациям NIST~\cite{nist-sp800-57}.}
\label{l9:keysizes}
\end{table}

Разрыв в таблице~\ref{l9:keysizes} растёт с уровнем стойкости: длина ключа RSA и DH в $\Znx{p}$
растёт быстрее линейной из-за субэкспоненциальных алгоритмов (решето числового поля),
а длина ключа ECC "--- линейно, вдвое больше стойкости.

\subsection{Почему не работает исчисление индексов}

Метод исчисления индексов в $\Znx{p}$ (лекция~4) опирается на два свойства:
\begin{enumerate}
  \item элементы $\Znx{p}$ можно \emph{поднять} в $\ZZ$ (вычет $\to$ число из $[1, p - 1]$),
    где есть разложение на простые множители;
  \item доля \emph{гладких} чисел (раскладывающихся по небольшой факторной базе) достаточно
    велика, и случайные степени $g^k$ гладки с заметной вероятностью.
\end{enumerate}
Для точек $E(\FF_p)$ естественный аналог "--- подъём на кривую над $\QQ$ "--- не даёт ни того,
ни другого: у рациональных точек нет разложения «на простые точки», ранг групп $E(\QQ)$ мал, а
\emph{высота} (размер координат) кратных точек растёт квадратично: координаты $kP$ имеют
порядка $k^2$ цифр. Попытки построить такой метод (например, «xedni calculus» Сильвермана)
оказались не быстрее перебора. Поэтому для «общих» кривых над простыми полями лучшие известные
алгоритмы "--- универсальные, сложности $\bigO(\sqrt q)$.

\subsection{Слабые кривые}

Есть специальные классы кривых, для которых ECDLP проще; стандарты их явно исключают.
\begin{itemize}
  \item \emph{Аномальные} кривые, $\#E(\FF_p) = p$ (след Фробениуса $t = 1$): ECDLP решается за
    полиномиальное время с помощью $p$-адического подъёма (Семаев, Сатох"--~Араки, Смарт,
    1998--1999, \cite{smart-anomalous}).
  \item \emph{Суперсингулярные} кривые (например, $\#E(\FF_p) = p + 1$) и вообще кривые
    с малой \emph{степенью вложения} $k$ ($q \mid p^k - 1$): атака MOV (Менезес"--~Окамото"--~Ванстоун,
    \cite{mov}) с помощью спаривания Вейля переводит ECDLP в $\langle P \rangle$ в задачу DLP
    в $\FF_{p^k}^*$, где работает субэкспоненциальное исчисление индексов. Для суперсингулярных
    кривых над $\FF_p$, $p > 3$, $k = 2$.
  \item Составной порядок точки без большого простого делителя "--- алгоритм Полига"--~Хеллмана.
\end{itemize}
Отсюда требования к кривой: $\#E(\FF_p) = hq$ с большим простым $q$ и малым $h$,
$\#E(\FF_p) \ne p$, $q \nmid p^k - 1$ при малых $k$ (на практике при $k \le 20$ или больше).

\checkyourself{Почему кривая $y^2 = x^3 + 2x + 2$ над $\FF_{17}$ «хороша» с точки зрения
алгоритма Полига"--~Хеллмана, а кривая с $\#E(\FF_p) = 2^{10}$ "--- нет? Сколько групповых
операций потребует алгоритм Полига"--~Хеллмана для точки порядка $2^{10}$?}

\section{Протоколы на эллиптических кривых}\label{l9:protocols}

\emph{Открытые параметры} (домен): простое $p$, кривая $E\colon y^2 = x^3 + ax + b$ над
$\FF_p$, точка $P \in E(\FF_p)$ простого порядка $q$ и кофактор $h = \#E(\FF_p)/q$.

\subsection{Протокол Диффи"--~Хеллмана (ECDH)}

\begin{center}
\begin{tikzpicture}[x=1cm,y=1cm]
  \node at (5,1.4) {$E/\FF_p,\quad P,\quad q = \ord P$};
  \node (A) at (0,0) {\personicon{caBlue}{Алиса}};
  \node (B) at (10,0) {\personicon{caRed}{Боб}};
  \draw[-{Stealth[length=3mm]},thick] (1,0.5) -- (9,0.5) node[midway,above] {$A = aP$};
  \draw[{Stealth[length=3mm]}-,thick] (1,-0.1) -- (9,-0.1) node[midway,below] {$B = bP$};
  \node[font=\small,align=left,anchor=north] at (0,-0.9)
    {$a \in [1, q - 1]$\\ $A \gets aP$\\ $K \gets aB$};
  \node[font=\small,align=left,anchor=north] at (10,-0.9)
    {$b \in [1, q - 1]$\\ $B \gets bP$\\ $K \gets bA$};
  \node (E) at (5,-1.9) {\personicon{caGray}{Ева}};
  \node[font=\small,anchor=north] at (5,-2.5) {знает: $E, P, A, B$};
\end{tikzpicture}
\end{center}

\begin{enumerate}
  \item Алиса выбирает секретное $a \in [1, q - 1]$ и отправляет $A = aP$;
    Боб выбирает $b$ и отправляет $B = bP$.
  \item Алиса вычисляет $K = aB$, Боб "--- $K = bA$. Так как $aB = a(bP) = (ab)P = b(aP) = bA$,
    ключи совпадают. Общим секретом обычно служит $x$-координата $x(K)$ (из неё ключ шифрования
    получают хеш-функцией).
\end{enumerate}

Ева знает $P$, $A$, $B$ и должна вычислить $abP$ "--- это \emph{задача Диффи"--~Хеллмана} на
кривой (ECDHP). Если Ева умеет решать ECDLP, то она найдёт $a = \log_P A$ и $K = aB$; обратное
(ECDHP $\Rightarrow$ ECDLP) не доказано, но лучших атак на ECDHP для хороших кривых не известно.

\begin{example}\label{l9:ex-ecdh}
$E\colon y^2 = x^3 + 2x + 2$ над $\FF_{17}$, $P = (5, 1)$, $q = 19$. Алиса: $a = 3$,
$A = 3P = (10, 6)$; Боб: $b = 10$, $B = 10P = (7, 11)$ (таблица в
примере~\ref{l9:ex-trivial}). Алиса: $2B = 20P = P = (5, 1)$, $K = 3B = P + B = (13, 10)$.
Боб: $2A = (16, 13)$, $4A = (0, 11)$, $5A = 4A + A = (3, 16)$, $K = 10A = 2 \cdot 5A = (13, 10)$.
Действительно, $K = 30P = 11P = (13, 10)$, общий секрет $x(K) = 13$.
\end{example}

\taskexample{ecdlp-ecdh}

\begin{remark}[проверка полученных точек]
Получив точку $B$, Алиса должна проверить, что $B \ne \ECpt$, что $B$ лежит на $E$ и что
$qB = \ECpt$. Иначе противник может подсунуть точку малого порядка (или точку на другой кривой
с теми же формулами сложения) и узнать $a \bmod \ord B$ перебором; несколько таких запросов
и китайская теорема об остатках дают весь ключ $a$ (атаки малой подгруппы и «неверной кривой»).
Секретные ключи, кратные $q$, дают $A = \ECpt$ и общий ключ $\ECpt$, известный всем.
\end{remark}

\begin{remark}[сжатие точки]
Вместо $(x, y)$ достаточно передать $x$ и один бит $y \bmod 2$: по $x$ находим
$y^2 = x^3 + ax + b$ и извлекаем квадратный корень (лекция~7); из двух корней $\pm y$ ровно один
чётный.
\end{remark}

\subsection{Шифрование Эль-Гамаля (ECElGamal)}

Открытый ключ Алисы $A = aP$, секретный "--- $a$. Сообщение "--- точка $M \in E(\FF_p)$.
\begin{itemize}
  \item \emph{Шифрование} (Боб): случайное $k \in [1, q - 1]$, шифртекст
    $(C_1, C_2) = (kP,\ M + kA)$.
  \item \emph{Расшифрование} (Алиса): $M = C_2 - aC_1$.
\end{itemize}
Корректность: $C_2 - aC_1 = M + k(aP) - a(kP) = M$. Пример на нашей кривой: $a = 3$,
$A = (10, 6)$; $M = 4P = (3, 1)$, $k = 2$: $C_1 = 2P = (6, 3)$, $C_2 = M + 2A = 4P + 6P = 10P =
(7, 11)$; Алиса вычисляет $aC_1 = 6P = (16, 13)$ и $C_2 - aC_1 = 4P = (3, 1) = M$. Недостатки:
сообщение нужно закодировать точкой, а шифртекст вчетверо длиннее открытого текста; на
практике используют гибридные схемы (ECIES): ECDH даёт ключ симметричного шифра.

\subsection{Подпись ECDSA}

Стандартная схема подписи на кривых~\cite{ecdsa} (аналог DSA). Ключи: $a$ и $A = aP$,
$h = H(m) \bmod q$ "--- хеш сообщения.
\begin{itemize}
  \item \emph{Подпись}: случайное $k \in [1, q - 1]$, $r = x(kP) \bmod q$,
    $s = k^{-1}(h + ar) \bmod q$ (если $r = 0$ или $s = 0$, выбрать другое $k$); подпись $(r, s)$.
  \item \emph{Проверка}: $u_1 = h s^{-1}$, $u_2 = r s^{-1} \bmod q$; подпись верна, если
    $x(u_1 P + u_2 A) \bmod q = r$.
\end{itemize}
Корректность: $u_1 P + u_2 A = s^{-1}(h + ar) P = kP$. Число $k$ должно быть секретным и
\emph{новым} для каждой подписи: по двум подписям с одним $k$ ключ $a$ находится решением
линейной системы (именно так в 2010 г. был извлечён ключ подписи игровой консоли PlayStation~3).

\checkyourself{Алиса и Боб используют ECDH с точкой $P$ порядка $q$, а Ева подменила $B$ на
точку $B'$ порядка $3$. Что Ева узнает об $a$, если Алиса затем использует ключ $K = aB'$?}

\section{\texorpdfstring{$\rho$}{rho}-метод Полларда для ECDLP}\label{l9:rho}

Напомним (лекция~3): если $f\colon S \to S$ "--- «случайное» отображение конечного множества
$|S| = N$, то последовательность $R_{j+1} = f(R_j)$ зацикливается (форма буквы $\rho$), и
первое совпадение $R_i = R_j$ ожидается после $\approx \sqrt{\pi N / 2}$ шагов. Метод Флойда
находит совпадение вида $R_i = R_{2i}$, храня только две точки: «черепаху» $R_i$ и «зайца»
$R_{2i}$.

\subsection{Псевдослучайное блуждание на кривой}

Пусть $Q = nP$, $N = \ord P$. Разобьём $E(\FF_p)$ на три множества примерно равного размера,
например по $x$-координате:
\[
  S_0 = \{R : x(R) \equiv 0\} \cup \{\ECpt\}, \qquad S_1 = \{R : x(R) \equiv 1\}, \qquad
  S_2 = \{R : x(R) \equiv 2\} \pmod 3,
\]
и определим отображение $f$ и пересчёт коэффициентов:
\[
  f(R) = \begin{cases} R + P, & R \in S_0,\\ 2R, & R \in S_1,\\ R + Q, & R \in S_2,\end{cases}
  \qquad
  (a, b) \mapsto \begin{cases} (a + 1, b),\\ (2a, 2b),\\ (a, b + 1) \end{cases} \pmod N.
\]
Если $R = aP + bQ$, то $f(R) = a'P + b'Q$ для новых $(a', b')$, поэтому все точки блуждания
известны в виде $R_j = a_j P + b_j Q$, хотя $n$ неизвестно. Начинаем с $R_0 = P$
($a_0 = 1$, $b_0 = 0$) или с $R_0 = a_0 P + b_0 Q$ при случайных $a_0, b_0$.

\begin{lemma}\label{l9:rho-congr}
Пусть $R_i = a_i P + b_i Q = a_j P + b_j Q = R_j$. Тогда
\[
  (b_i - b_j)\, n \equiv a_j - a_i \pmod N.
\]
Если $d = \gcd(b_i - b_j, N)$, то сравнение имеет ровно $d$ решений по модулю $N$, и $n$ "---
одно из них. В частности, при простом $N$ и $b_i \not\equiv b_j$ решение единственно:
$n \equiv (a_j - a_i)(b_i - b_j)^{-1} \pmod N$.
\end{lemma}

\begin{proof}
Подставляя $Q = nP$: $(a_i + b_i n) P = (a_j + b_j n) P$, т.\,е. $N \mid (a_i + b_i n) - (a_j + b_j n)$.
Линейное сравнение $\beta x \equiv \alpha \pmod N$ с $d = \gcd(\beta, N)$ разрешимо тогда и
только тогда, когда $d \mid \alpha$ (здесь разрешимо, так как $x = n$ "--- решение), и тогда
равносильно сравнению $\frac{\beta}{d} x \equiv \frac{\alpha}{d} \pmod{\frac Nd}$ с обратимым
коэффициентом; его решение $x_0$ единственно по модулю $N/d$ и даёт $d$ решений
$x_0 + k\frac Nd$, $0 \le k < d$, по модулю $N$.
\end{proof}

Если $d > 1$ невелико, перебираем $d$ кандидатов, проверяя $x P = Q$. Если же
$b_i \equiv b_j \pmod N$ (тогда и $a_i \equiv a_j$), то коллизия бесполезна: сравнение
$0 \cdot n \equiv 0$ ничего не даёт, и блуждание перезапускают с другой начальной точкой.
Для простого $N$ это происходит с вероятностью около $1/N$.

\begin{algorithm}[$\rho$-метод Полларда для ECDLP]\label{l9:alg-rho}\leavevmode
\begin{algorithmic}[1]
\Require $P$ порядка $N$, $Q \in \langle P \rangle$.
\Ensure $n$: $Q = nP$.
\State $(T, a_T, b_T) \gets (R_0, a_0, b_0)$, \quad $(H, a_H, b_H) \gets (R_0, a_0, b_0)$
\Repeat
  \State $(T, a_T, b_T) \gets f(T, a_T, b_T)$ \Comment{черепаха: один шаг}
  \State $(H, a_H, b_H) \gets f(f(H, a_H, b_H))$ \Comment{заяц: два шага}
\Until{$T = H$}
\If{$b_T \equiv b_H \pmod N$} \State перезапуск с другими $a_0, b_0$ \EndIf
\State решить $(b_T - b_H)\, n \equiv a_H - a_T \pmod N$ (лемма~\ref{l9:rho-congr}), проверить
  кандидатов
\State \Return $n$
\end{algorithmic}
\end{algorithm}

\begin{theorem}[эвристическая оценка]\label{l9:rho-cost}
Если отображение $f$ ведёт себя как случайное, то алгоритм~\ref{l9:alg-rho} решает ECDLP за
$\bigO(\sqrt N)$ групповых операций (в среднем) и $\bigO(1)$ памяти.
\end{theorem}

\begin{proof}[Обоснование]
Последовательность $R_j$ лежит в $\langle P \rangle$ ($N$ элементов). По теореме о коллизиях
(лекция~3) для случайного отображения ожидаемые длины предпериода $t$ и цикла $c$ равны
$\bigO(\sqrt N)$ (каждая $\approx \sqrt{\pi N / 8}$). Метод Флойда находит первое $i \ge t$,
кратное $c$, т.\,е. $i < t + c$, и делает $3i = \bigO(\sqrt N)$ вычислений $f$, каждое "---
одна групповая операция. Решение сравнения стоит $\bigO(\log N)$ операций в $\Zn N$.
Строго доказать «случайность» конкретного $f$ не удаётся, поэтому оценка эвристическая; на
практике она хорошо подтверждается (разбиение на $3$ множества даёт блуждание чуть хуже
случайного; Э.~Теске предложила разбиение на $r \approx 20$ множеств с $R \mapsto R + M_s$,
$M_s = \alpha_s P + \beta_s Q$).
\end{proof}

\begin{example}\label{l9:ex-rho}
$E\colon y^2 = x^3 + 2x + 2$ над $\FF_{17}$, $P = (5, 1)$, $N = 19$, $Q = (16, 4)$. Блуждание
с $R_0 = P$ ($s = x(R_j) \bmod 3$ "--- номер множества):
\[
\begin{array}{r|c|c|r|r}
  j & R_j & s & a_j & b_j \\ \hline
  0 & (5, 1)   & 2 & 1  & 0  \\
  1 & (9, 1)   & 0 & 1  & 1  \\
  2 & (3, 16)  & 0 & 2  & 1  \\
  3 & (10, 11) & 1 & 3  & 1  \\
  4 & (16, 4)  & 1 & 6  & 2  \\
  5 & (0, 6)   & 0 & 12 & 4  \\
  6 & (13, 7)  & 1 & 13 & 4  \\
  7 & (10, 11) & 1 & 7  & 8  \\
  8 & (16, 4)  & 1 & 14 & 16
\end{array}
\]
Предпериод $t = 3$, цикл $c = 4$ ($R_3 = R_7$). Флойд: $R_1 \ne R_2$, $R_2 \ne R_4$,
$R_3 \ne R_6$, $R_4 = R_8 = (16, 4)$. Отсюда $6P + 2Q = 14P + 16Q$,
\[
  (2 - 16)\, n \equiv 14 - 6, \qquad 5n \equiv 8 \pmod{19}, \qquad
  5^{-1} \equiv 4, \qquad n \equiv 32 \equiv 13 \pmod{19}.
\]
Проверка по таблице примера~\ref{l9:ex-trivial}: $13P = (16, 4) = Q$.
\end{example}

\taskexample{ecdlp-rho}

\taskexample{ecdlp-rho-composite}

\subsection{Ускорения}

\begin{remark}[отображение $R \mapsto -R$]
Точки $R$ и $-R$ имеют одинаковую $x$-координату, и их легко отождествить: выберем
представителя класса $\{R, -R\}$ (например, с меньшим $y$) и будем вести блуждание на
множестве классов, которых $\approx N/2$. Если $\pm R_i = \pm R_j$, то
$a_i + b_i n \equiv \pm (a_j + b_j n) \pmod N$ "--- снова линейное сравнение. Ожидаемое число
шагов уменьшается в $\sqrt 2$ раз (М.~Винер, Р.~Цуккерато~\cite{wiener-zuccherato}); на
практике нужно бороться с короткими «бесплодными циклами» блуждания.
\end{remark}

\begin{remark}[выделенные точки и распараллеливание]
Метод выделенных точек (лекция~3, ван~Ооршот и Винер~\cite{van-oorschot-wiener}): точку
назовём \emph{выделенной}, если, например, $\theta$ младших битов её $x$-координаты равны нулю.
Каждый из $m$ процессоров ведёт своё блуждание из случайной точки $a_0 P + b_0 Q$ и отправляет
на центральный сервер только выделенные точки вместе с $(a, b)$. Совпадение двух выделенных
точек (с разными $b$) даёт сравнение леммы~\ref{l9:rho-congr}. Время уменьшается в $m$ раз:
$\bigO(\sqrt N / m)$ на процессор. Так в 2009 г. был решён ECDLP на кривой над $112$-битовым
простым полем (кластер из $200$ игровых консолей, несколько месяцев). Для $N \approx 2^{256}$
нужно около $2^{128}$ операций "--- недостижимо.
\end{remark}

\checkyourself{В примере~\ref{l9:ex-rho} коллизия $R_3 = R_7$ встретилась раньше, чем её нашёл
метод Флойда. Какое сравнение даёт она? Получается ли тот же ответ?}

\checkyourself{Почему в $\rho$-методе нельзя взять $f(R) = 2R$ для всех $R$? А $f(R) = R + P$?}

\checkyourself{Пусть $\ord P = N = 2 \cdot 13$ и коллизия дала сравнение $4n \equiv 10
\pmod{26}$. Сколько кандидатов нужно проверить? Как лучше решать ECDLP в группе такого порядка?}

\checkyourselfsection

\begin{furtherreading}
  \readitem[разд.~6.3--6.4]{hps}{алгоритм «удвоение"--~сложение», тернарное разложение, ECDLP,
    ECDH и ECElGamal; разд.~5.5 "--- $\rho$-метод}
  \readitem[гл.~5--6]{washington}{ECDLP (в том числе MOV и аномальные кривые) и протоколы ECC}
  \readitem[гл.~11, 14]{galbraith}{эффективное возведение точки в степень (NAF, окна, лесенка);
    $\rho$-метод и его варианты, отображение $R \mapsto -R$}
  \readitem[разд.~14.6--14.7]{hac}{алгоритмы возведения в степень и знаковые разложения}
  \readitem[гл.~9]{paar-pelzl}{криптосистемы на эллиптических кривых, ECDH, длины ключей}
  \readitem{ecdsa}{подробное описание ECDSA и требований к параметрам}
  \readitem{mov}{сведение ECDLP к DLP в конечном поле с помощью спаривания Вейля}
  \readitem[лекции об ECC]{paar-lectures}{видеолекции К.~Паара по эллиптическим кривым}
  \readitem{mit-18783}{курс MIT по эллиптическим кривым (конспекты лекций)}
\end{furtherreading}
\end{lecture}
